\begin{center}\rule{0.5\linewidth}{0.5pt}\end{center}

\textbf{Alexander Winkler} ist wissenschaftlicher Angestellter bei
digiS, dem Forschungs- und Kompetenzzentrum Digitalisierung Berlin
(https://www.digis-berlin.de/), wo er sich aktuell insbesondere mit
Metadatenqualität, Linked Data, Datennutzung und KI im GLAM-Sektor
beschäftigt. Er hat Latinistik und Italianstik,
Renaissancewissenschaften und Bibliotheks- und
Informationswissenschaften in München, Pisa, Warwick und Berlin studiert
und über ein neulateinisch-italianistisches Thema promoviert. ORCID:
\url{https://orcid.org/0000-0002-9145-7238} \textbar{} Mastodon:
awinkler@openbiblio.social

\textbf{Jens Bemme} studierte Verkehrswirtschaft und
Lateinamerikastudien. Als Mitarbeiter der SLUB Dresden begleitet er
landeskundliche Open-Citizen-Science-Initiativen insbesondere mit
Werkzeugen, Daten, Methoden und Gemeinschaften der Wikimedia-Bewegung.
Tafellieder, Tanzkarten und die gemeinfreien Publikationen historischer
Geschichtsvereine stehen derzeit im Fokus. ORCID:
\url{https://orcid.org/0000-0001-6860-0924} \textbar{} Mastodon:
JensB@openbiblio.social
